% This section summarizes the verified full closed smoothed-power family.
% Formal scope and declarations are in Appendix~\ref{app:transport-provenance}.
% The generic convex spectral-penalty criterion in spectral_penalty_transport.tex
% has NOT been established, and is deliberately not imported here.
\section{Spectral-penalty transport: a common variational problem}
\label{sec:unified-transport}

Let $P,Q$ be orthogonal projections on a complex Hilbert space $E$.
The \emph{transport fiber} is
\begin{equation}
 \mathcal T(P,Q)=\{W\in\mathcal U(E):WPW^*=Q\}.
 \label{eq:transport-fiber}
\end{equation}
For a unitary $W$, the singular values of $I-W$ measure its
\emph{ambient displacement}.  The fiber condition fixes the image of
$\operatorname{ran}P$, while allowing $W$ to act freely within that subspace
and its orthogonal complement.  A principal-plane construction supplies
feasible transports, but it does not by itself exclude lower-cost
transformations mixing different principal planes
\citep{davis1958separation,davis1970rotation,bolotnikov2001minimal}.

We consider the closed two-parameter family of nonnegative penalties
\begin{equation}
 \rho_{p,\tau}(s)
   =(\tau^2+s^2)^{p/2}-\tau^p,
 \qquad 1\le p\le2,\quad \tau\ge0,\quad s\ge0,
 \label{eq:family-rho}
\end{equation}
where $\rho_{p,0}(s)=s^p$.  Define
\begin{equation}
 \mathcal F_{p,\tau}(W)=
    \operatorname{Tr}\rho_{p,\tau}(|I-W|),\qquad
 \phi_{p,\tau}(t)=(\tau^2+2-2\cos t)^{p/2}-\tau^p,
 \label{eq:family-objective}
\end{equation}
and, for $0\le\theta\le\pi/2$,
\begin{equation}
 m_{p,\tau}(\theta)=
 \min_{0\le\delta\le\theta}\bigl\{
      \phi_{p,\tau}(\theta+\delta)
     +\phi_{p,\tau}(\theta-\delta)\bigr\}.
 \label{eq:family-planar-minimum}
\end{equation}
The scalar phase $\delta$ adjusts the two eigenphases of a
principal-plane transporter.  Equation~\eqref{eq:family-planar-minimum}
by itself optimizes only that planar family; the next theorem concerns
\emph{every} ambient transporter.

\begin{theorem}[Exact unrestricted complex transport for the closed family]
\label{thm:family-finite}
Let $P,Q$ be equal-rank orthogonal projections on a finite-dimensional
complex Hilbert space, and let
$\theta_1,\ldots,\theta_r\in[0,\pi/2]$ be their principal angles,
with multiplicity.  For every $1\le p\le2$ and $\tau\ge0$,
\begin{equation}
 \min_{W\in\mathcal T(P,Q)}\mathcal F_{p,\tau}(W)
   =\sum_{j=1}^r m_{p,\tau}(\theta_j).
 \label{eq:family-finite-minimum}
\end{equation}
The minimum is attained by a unitary using independently minimizing
phases in a cosine--sine decomposition and fixing the unused common
directions.  The lower bound holds for all ambient unitaries satisfying
$WPW^*=Q$, without a principal-plane preservation assumption.
\end{theorem}

\begin{proof}[Structure of the formalized argument]
For $\tau>0$, compactness of the fiber gives a minimizing unitary.
First variation along the unitary stabilizer of $P$ forces an
anticommutation identity between the spectral gradient and the source
reflection $2P-I$.  Second variation excludes overlapping large-phase
spectral subspaces.  The resulting decomposition consists of paired
transport planes and a zero-gradient reflection sector.  Each plane
has cost at least~\eqref{eq:family-planar-minimum}; the reflection
sector is bounded using the nuclear projection-distance inequality
and the convexity comparison
\[
 m_{p,\tau}(\theta)
   \le\rho_{p,\tau}(2\sin\theta)
   \le\rho_{p,\tau}(2)\sin\theta.
\]
An independent-phase construction attains the sum of these lower
bounds.  At $\tau=0$, continuity of the finite-dimensional spectral
cost and its compact minimum, together with an explicit limiting
construction, yields the closed boundary.  These steps are certified
by the finite-dimensional spectral-decomposition and boundary
theorems identified in Appendix~\ref{app:transport-provenance}.
\end{proof}

\subsection{A sharp phase-transition surface}
\label{sec:family-phase}

Set
\begin{equation}
 C_{p,\tau}=
 \frac{\tau^2+2-\sqrt{(\tau^2+2)^2-4p(2-p)}}{2p}.
 \label{eq:family-critical-cosine}
\end{equation}
This threshold belongs to $[0,1]$ for the indicated parameters.
For $\tau>0$ and $0<\theta<\pi/2$, the unique minimizing phase
in~\eqref{eq:family-planar-minimum} is $\delta=0$ when
$\cos\theta\ge C_{p,\tau}$.  When $\cos\theta<C_{p,\tau}$,
there is a unique positive minimizing phase, characterized by
\[
 \phi_{p,\tau}'(\theta+\delta)
   =\phi_{p,\tau}'(\theta-\delta).
\]
The same canonical-minimum criterion extends to the closed parameter
range and all $\theta\in[0,\pi/2]$, with the endpoint conventions
spelled out below.

\begin{theorem}[Ambient canonical-optimality criterion]
\label{thm:family-canonical}
Under the hypotheses of \cref{thm:family-finite}, a completed canonical
direct rotation is a global minimizer of $\mathcal F_{p,\tau}$ over
$\mathcal T(P,Q)$ if and only if
\begin{equation}
  \cos\theta_j\ge C_{p,\tau}\qquad\text{for every }j.
  \label{eq:family-canonical-criterion}
\end{equation}
\end{theorem}

The threshold comes from the scalar curvature identity
\[
 \phi_{p,\tau}''(t)
  =p(\tau^2+2-2\cos t)^{p/2-2}
     \bigl[p-2+(\tau^2+2)\cos t-p\cos^2t\bigr]
\]
where the expression is evaluated away from points at which the
unsmoothed derivative is singular.  The formally proved
single-crossing argument converts the curvature change into a global
phase-minimum statement; the ambient equivalence then follows from
\cref{thm:family-finite}.  Its boundary cases include
$C_{p,0}=(2-p)/p$, the nuclear corner $C_{1,0}=1$, and
$C_{2,\tau}=0$.  At $p=2$, the phase cost is independent of $\tau$
and is flat at a right angle; \cref{thm:family-canonical} asserts
minimality there, without asserting uniqueness.

The two-parameter theorem supplies minimum values and an attaining
construction.  Complete classifications of \emph{all} minimizers and
local-to-global conclusions have separately been proved for the strict
unsmoothed Schatten regime, the nuclear endpoint, and the original
Charbonnier penalty; the family theorem alone does not extend those
classifications to every interior pair $(p,\tau)$.

\subsection{Operator ideals and infinite-dimensional attainment}
\label{sec:family-infinite}

For a complete, possibly nonseparable complex Hilbert space, the trace
in~\eqref{eq:family-objective} is the extended positive trace.
In particular, a nonempty transport fiber can have infinite cost for
every transporter.  Define the crossed defects
\[
 D_{P\bar Q}=\operatorname{ran}P\cap\ker Q,
 \qquad D_{\bar P Q}=\ker P\cap\operatorname{ran}Q.
\]
For $\Delta=P-Q$, the necessary ideal depends on smoothing:
\begin{equation}
 q=\begin{cases}
       p,&\tau=0,\\
       2,&\tau>0.
    \end{cases}
 \label{eq:family-ideal-exponent}
\end{equation}
The change occurs because $\rho_{p,0}(s)=s^p$, whereas for
$\tau>0$ one has
$\rho_{p,\tau}(s)\sim \tfrac p2\tau^{p-2}s^2$
as $s\downarrow0$.

\begin{theorem}[Finite-cost feasibility and exact Hilbert-space minimum]
\label{thm:family-infinite}
Let $P,Q$ be orthogonal projections on a complete complex Hilbert
space, $1\le p\le2$, and $\tau\ge0$.  A transporter
$W\in\mathcal T(P,Q)$ with $\mathcal F_{p,\tau}(W)<\infty$
exists if and only if
\begin{equation}
 \Delta=P-Q\in S_q,
 \qquad \dim D_{P\bar Q}=\dim D_{\bar P Q},
 \label{eq:family-ideal-and-defect}
\end{equation}
where $q$ is given in~\eqref{eq:family-ideal-exponent}.
Under these conditions the infimum is attained and
\begin{equation}
 \min_{W\in\mathcal T(P,Q)}\mathcal F_{p,\tau}(W)
   =\sum_j m_{p,\tau}(\theta_j)<\infty.
 \label{eq:family-infinite-minimum}
\end{equation}
Here $(\theta_j)$ lists the positive principal angles determined by
the nonzero singular values of $(I-Q)P$, counted with multiplicity;
zero angles contribute zero.  The list is countable under the ideal
hypothesis, even when the ambient Hilbert space is nonseparable.
\end{theorem}

The lower bound applies to arbitrary ambient transporters through
finite-section approximation and passage to the spectral-angle sum.
Existence uses a balanced-defect completion of the direct rotation and
compatible phase corrections; the nuclear corner requires a distinct
signed-root construction.  The relative-index obstruction is classical
\citep{avron1994index}, while the attained value for this spectral
penalty family follows from the formally verified transport theorem.
The existence theorem does not assert an all-optimizer classification
for arbitrary $(p,\tau)$.

\subsection{Real-field compatibility and the nuclear endpoint}
\label{sec:family-real}

The complex phase selections in \cref{thm:family-finite} need not be
realizable by independent real principal-plane rotations.  For
finite-dimensional real equal-rank subspaces and unsmoothed Schatten
powers $1\le p<2$, the separately proved unrestricted real minimum is
an attained weighted matching of the principal-angle labels.
At $p=1$, writing $J=\{i:\theta_i>0\}$, a matching $M$ consists of
disjoint unordered pairs from $J$, and its cost is
\begin{equation}
 \sum_{\{i,j\}\in M}4\sin\frac{\theta_i+\theta_j}{2}
 +\sum_{i\in J\setminus\cup M}
          \min\{4\sin(\theta_i/2),2\}.
 \label{eq:family-real-matching}
\end{equation}
The real nuclear minimum is the smallest such cost over all matchings.
The complex counterpart is $2\sum_j\sin\theta_j$.
The two values are equal precisely when every strictly acute positive
principal angle has even multiplicity; zero and right angles are
unrestricted.  \Cref{sec:nuclear-endpoint} gives the finite-dimensional
theorems and the sharp condition for equality with the real canonical
cost.

In arbitrary complete real Hilbert spaces, for trace-class
$P-Q$ and balanced crossed defects, the complex nuclear lower bound
$\|P-Q\|_1$ is attained by a real transporter exactly when each
strictly acute positive sine-eigenspace has even \emph{source-side}
multiplicity.  All equality attainers admit the checked signed-root
certificate.  This extends the explanation of the finite parity rule
to compact spectral blocks; it does not give an unrestricted real
infinite-rank matching formula when parity fails.

Consequently, the complex closed-family theorem, the real finite
matching result, and the real nuclear compatibility criterion form a
field-sensitive theory with precisely different scopes.  In particular,
no unrestricted real analogue of \cref{thm:family-finite} for positive
smoothing and arbitrary $p$ is claimed here.
